\section{经济犯罪 - 金融诈骗罪、危害税收征管罪、侵犯知识产权罪、扰乱市场秩序罪}
\begin{enumerate}
  \item 金融诈骗罪。
        \begin{enumerate}
          \item 贷款诈骗。\zd{虚假担保问题}。
                \begin{enumerate}
                  \item \al{甲骗到乙的汽车，质押给银行贷到款跑路}。
                  \item \al{甲借用到乙的汽车，质押给银行贷到款跑路}。
                  \item \al{甲骗乙做担保，乙用自己的汽车质押给银行，甲贷到款跑路}。\zd{手段与目的，牵连犯}。
                \end{enumerate}
          \item 保险诈骗。采取虚构保险标的、保险事故或者制造保险事故等方法，骗取保险金。
                \begin{enumerate}
                  \item 主体特殊。投保人、被保险人、受益人。\al{汽修厂老板甲多次向保险公司夸大汽车损毁程度}。
                  \item 侵害对象。\zd{仅包括商业保险，不包括社会保险（三险一金）}
                  \item 罪数问题。\al{杀人骗保，申请理赔着手}。
                \end{enumerate}
          \item \bt{信用卡诈骗}。最难。
                \begin{enumerate}
                  \item 行为对象。信用卡扩大解释 = 银行卡。
                  \item 行为方式。对人、对机器都定信用卡诈骗。
                        \begin{enumerate}
                          \item 使用伪造（伪造金融票证）\hz 骗领（妨害信用卡管理）来的信用卡。\zd{使用假卡，牵连犯}。
                          \item 冒\zd{用}他人信用卡。不用无罪，用了一律定信用卡诈骗；\zd{抢劫/盗窃 + 使用 = 抢劫/盗窃}。
                                \begin{itemize}
                                  \item 考试。判断非法取得行为。\al{甲帮好友乙打扫房间，捡到一张信用卡}。
                                  \item 放到共犯平台 + 认识错误。\al{甲盗窃丙信用卡，骗乙说捡来的，乙商场消费 3W}。
                                  \item 用包括自己用，给别人用，甲既是教唆，同时也是帮助。
                                \end{itemize}
                          \item 恶意透支。\zd{合法的持卡人；两次催收还不还}。
                        \end{enumerate}
                  \item 非法获取银行卡信息资料再用，一律定信用卡诈骗。
                  \item 微信/支付宝信息资料。盗划、盗刷 \yjt 盗窃罪。
                  \item 蚂蚁花呗/借呗。本质是消费信贷，垫付/放贷资金 \yjt 贷款诈骗。
                \end{enumerate}
          \item 集资类犯罪。
                \begin{enumerate}
                  \item 非法吸收公众存款。没有非法占有目的，不想还 \yjt 集资诈骗。
                  \item 集资诈骗。有非法占有目的。吴英案。
                  \item 擅自发行股票、公司、企业债券。
                  \item 非法经营。
                \end{enumerate}
        \end{enumerate}
  \item 危害税收征管罪。
        \begin{enumerate}
          \item 逃税。
                \begin{enumerate}
                  \item 有第一款行为，经\zd{税务机关}依法下达追缴通知后，补缴应纳税款，缴纳滞纳金，已受行政处罚的，不予追究刑事责任；
                        但是，五年内因逃避缴纳税款受过刑事处罚或者被税务机关给予二次以上行政处罚的除外。
                  \item 只有纳税人能享受，单位不行。
                  \item \zd{处罚阻却事由，独一份}。
                \end{enumerate}
          \item 骗取出口退税。
                \begin{enumerate}
                  \item \zd{想象竞合数罪并罚。独一份}。
                  \item \al{甲依法缴纳 100W 税，骗取出口退税方式骗回 120W}。逃税 100W + 骗税 20W。
                \end{enumerate}
          \item 虚开增值税专用发票。税收安全。
                \begin{enumerate}
                  \item 目的/行为 1：虚开目的/行为；目的/行为 2：Required(逃税的目的)/Opt(行为)
                  \item 目的 1 + \zd{目的 2} \yjt 行为 1（虚开）。
                  \item 认知错误。以为是专票其实是普票 \yjt 虚开发票罪。
                \end{enumerate}
          \item 虚开发票罪。普通发票。
        \end{enumerate}
  \item 侵犯知识产权罪。
        \begin{enumerate}
          \item 侵犯著作权。民法里画个小圈。
                \begin{enumerate}
                  \item 侵犯方式。发行、网络传播。
                  \item 发行。\zd{内容 + 载体转让}给不特定人或多数人。
                  \item 播放。\zd{内容展示}给不特定人或多数人，载体未转让。
                  \item \zd{线上发行、线下发行、线上播放构成}，线下播放不构成。
                \end{enumerate}
          \item 侵犯商业秘密。方式。不正当手段获取；程度要求。造成重大财产损失改为情节严重。情节犯。
          \item 为境外窃取、刺探、收买、非法提供商业秘密。商业间谍。
          \item 假冒注册商标罪。增加了对服务类（理发店）商标的保护。
        \end{enumerate}
  \item 扰乱市场秩序罪。
        \begin{enumerate}
          \item 非法经营罪。
                \begin{enumerate}
                  \item 烟草。专营。
                  \item 口袋罪：其他严重扰乱市场秩序的行为。
                        \begin{enumerate}
                          \item 经营非法出版物，\zd{有没有合法的著作权}；从事非法网络水军活动；疫情期间哄抬物价。
                          \item 擅自设立印刷厂、出版社、图书进口单位。
                          \item \zd{放高利贷：盈利目的；经常性；向社会不特定对象放}。
                        \end{enumerate}
                \end{enumerate}
          \item 强迫交易罪。强买强卖；强迫他人提供 \hz 接受服务。
          \item 提供虚假证明文件。\al{公司上市财物报告做假账}。牵连犯，择一重。
          \item 合同诈骗。经济合同。
          \item 组织、领导传销活动。没有参加；骗取财物 \yjt 间接目的，动机，不要求客观行为。
        \end{enumerate}
\end{enumerate}
